\chapter{Stochastic Trajectories \& Motor Variability}
\label{ch:stochastic_trajectories}

\begin{gomdraftnotice}
This chapter is under active development and contains only a structural outline
with preliminary material. Full exposition, worked examples, proofs, and exercises
will be added in future editions.
\end{gomdraftnotice}

\epigraph{%
  Perfection is inherently fragile. Biological optimization seeks the strongest balance between intent and reality.
}{}

\vspace{0.5cm}

\begin{laymansbox}{The Harder You Try, The More You Miss}{harder_you_try}
If you program a robot motor to apply 50 units of force, it applies exactly 50 units of force, every single time. Human muscles don't work like that. The harder you try to contract a muscle, the more "noisy" and jittery the actual force becomes. This is a massive problem for trying to swing a golf club as fast as humanly possible: the closer you get to your maximum physical effort, the more unpredictable your motion becomes. This chapter explains the math behind why smooth, controlled movements are more accurate than jerky, maximum-effort ones. We'll see that a truly "optimal" golf swing doesn't use 100\% of your strength, because dialing it back slightly eliminates a huge amount of physical chaos, making the swing actually repeatable.
\end{laymansbox}

% ══════════════════════════════════════════════
\section{Signal-Dependent Noise}
% ══════════════════════════════════════════════

When a roboticist specifies a torque of $50\text{ Nm}$, the motor delivers $50\text{ Nm}$ with a tiny, constant variance. Biological actuators do not work this way. Human muscles generate force by recruiting discrete motor units \cite{Henneman1957}. As higher forces are required, larger, more fast-twitch oriented motor units are recruited. This physiological reality manifests as a fundamental control challenge: \textbf{motor noise is proportional to the command signal} \cite{HarrisWolpert1998}.

\begin{definition}{Signal-Dependent Noise (SDN)}{sdn}
  A control input $\control(t)$ applied by a biological system is actually realized as a stochastic variable $\tilde{\control}(t)$:
  \begin{equation}
    \tilde{\control}(t) = \control(t) + \bm{W} \control(t) \cdot \bm{\epsilon}(t)
  \end{equation}
  where $\bm{W}$ is a scaling matrix representing the coefficient of variation, and $\bm{\epsilon}(t) \sim \mathcal{N}(0, I)$ is white noise. That is, the variance of the applied force grows squarely with the magnitude of the ordered force: $\text{Var}(\tilde{u}) = c u^2$.
\end{definition}

This simple reality crumbles deterministic optimal control. If you compute an optimal trajectory that requires maximum muscular exertion, the trajectory will mathematically possess minimum variance only if noise is constant. Under SDN, maximum exertion guarantees maximum chaotic noise.

% ══════════════════════════════════════════════
\section{Minimum Variance Theory of Human Movement}
% ══════════════════════════════════════════════

Why do humans move in smooth, bell-shaped velocity profiles rather than jerky "bang-bang" optimal control solutions when reaching for coffee?

One influential explanation is the \textbf{minimum jerk} theory of Flash and Hogan \cite{FlashHogan1985}, which posits that the nervous system minimizes $\int \dddot{x}^2 \dt$ (the integral of squared jerk). This kinematic criterion successfully predicts many features of reaching movements. An alternative, physically grounded explanation comes from the \textbf{Minimum Variance Theory} of Harris and Wolpert \cite{HarrisWolpert1998}, which derives smooth trajectories from the assumption of signal-dependent noise. Both theories remain active areas of research, and they are not mutually exclusive---minimum-jerk trajectories can emerge as a consequence of variance minimization under certain noise models.

\begin{keyidea}{Optimization for Repeatability}{min_variance}
  The human nervous system seeks to move in a way that minimizes the final endpoint variance evaluated across multiple trials. Under signal-dependent noise, this objective naturally selects for smooth trajectories with bounded control effort.
\end{keyidea}

Under SDN, the only way to minimize endpoint variance (ensure you actually grasp the coffee cup smoothly) is to minimize sudden spikes in required actuator forces, because large forces inject huge amounts of uncertainty into the dynamics. Smooth trajectories inherently possess lower peak torque commands than rapid accelerations and decelerations. Smoothness is not the explicitly programmed objective; it is an emergent property of attempting to be consistently accurate under the unavoidable reality of signal-dependent noise.

\begin{center}
\begin{tikzpicture}[scale=1.0]
  % Deterministic trajectory (top)
  \begin{scope}[yshift=2.2cm]
    \node[font=\bfseries, chapblue] at (-1.8, 1.2) {Deterministic OCP};

    % Reference trajectory
    \draw[very thick, chapblue] (0, 0) .. controls (1, 0.8) and (2, 0.6) .. (3, 0.1);
    \node[chapblue, font=\small] at (1.5, 1.0) {$\traj^*(t)$};

    % Point mass at end
    \filldraw[chapblue] (3, 0.1) circle (3pt);
    \node[chapblue, font=\small] at (3.5, 0.3) {Target};

    \node[gray, font=\small, align=center] at (1.5, -0.7) {Single solution\\(no robustness)};
  \end{scope}

  % Stochastic ensemble (bottom)
  \begin{scope}[yshift=-0.3cm]
    \node[font=\bfseries, trajgreen!80!black] at (-1.8, 1.2) {Stochastic OCP};

    % Reference trajectory
    \draw[very thick, trajgreen] (0, 0) .. controls (1, 0.8) and (2, 0.6) .. (3, 0.1);
    \node[trajgreen!80!black, font=\small] at (1.5, 1.0) {$\traj^*(t)$ (nominal)};

    % Perturbed trajectories (Monte Carlo samples)
    \foreach \offset in {-0.15, -0.10, -0.05, 0.05, 0.10, 0.15} {
      \draw[thick, accentred!40, opacity=0.5] (0, \offset) .. controls (1, 0.75 + \offset) and (2, 0.55 + \offset) .. (3, 0.05 + \offset);
    }

    % Terminal covariance ellipse
    \draw[dashed, accentred, line width=1.5pt] (3, 0.1) ellipse (0.35 and 0.15);
    \node[accentred, font=\small] at (3.8, 0.3) {$\mathbb{E}[\psi]$};

    % SDN annotation
    \node[red!70, font=\small, align=center] at (0.5, -0.5) {$\text{Var}(\tilde{u}) \propto u^2$};
    \node[gray, font=\small, align=center] at (1.5, -0.7) {Ensemble of trajectories\\under signal-dependent noise};
  \end{scope}

  % Dividing line
  \draw[thick, gray!40] (-2.3, 1.8) -- (4.3, 1.8);
\end{tikzpicture}
\end{center}

% ══════════════════════════════════════════════
\section{The Speed-Accuracy Tradeoff in the Golf Swing}
% ══════════════════════════════════════════════

Fitts's Law \cite{Fitts1954} states that the faster you try to move to a target, the less accurate you become. In golf, we see the supreme manifestation of this phenomenon. The deterministic optimal control solution (Chapter 6) tells the solver to maximize input torque to achieve theoretically maximal clubhead speed. 

However, if we introduce the Stochastic OCP (via Covariance Steering) from Chapter 6 and apply the realistic premise of SDN, the optimizer faces a stark tradeoff.

Applying maximum $u_1$ torque during the downswing injects enormous variance into the kinematic state. The unactuated club ($\theta_2$) begins whipping through its zero dynamics, but its exact phase progression is now highly uncertain due to the noise flooding the arms. 

To ensure the clubface squares up at the moving goal, the stochastic optimizer will \emph{deliberately back off} the peak requested torque. It trades a small reduction in peak terminal performance for a substantial reduction in angular variance at the target manifold constraint. The optimal biological swing is mathematically proven to be a sub-maximal physical exertion when signal-dependent noise is accounted for.

% ══════════════════════════════════════════════
\section{Summary}
% ══════════════════════════════════════════════

\begin{center}
\renewcommand{\arraystretch}{1.4}
\begin{tabular}{@{} l p{7.5cm} @{}}
  \toprule
  \textbf{Concept} & \textbf{Implication for Control} \\
  \midrule
  Signal-Dependent Noise (SDN) & Muscle noise variance proportional to command: $\text{Var}(\tilde{u}) = c u^2$. Requires bounded control effort for repeatability. \\
  Minimum Variance Theory & Explains smooth, controlled velocity profiles as solutions to endpoint error variance minimization under SDN. \\
  Stochastic OCP & Optimal control formulation that minimizes $\mathbb{E}[\phi(\state(t_f))] + \lambda \cdot \mathrm{Var}[\psi(\state(t_f))]$. \\
  Covariance Steering & Technique to propagate the covariance matrix $\bm{\Sigma}(t)$ alongside the nominal trajectory during optimization. \\
  Speed-Accuracy Tradeoff & Formal manifestation of Fitts's Law: higher peak forces inject noise proportional to force squared. \\
  \bottomrule
\end{tabular}
\end{center}

\vspace{0.3cm}


The extension of the deterministic Pontryagin framework to stochastic systems is developed rigorously in Agrachev and Sachkov \cite{AgrachevSachkov2004}. When the system is subject to signal-dependent noise (SDN), the traditional Hamiltonian maximization is supplemented by terms encoding risk: the optimizer must maximize expected performance while managing the variance of terminal conditions under noise propagation. This leads to \textbf{risk-sensitive control} problems where the cost includes an entropy-like term penalizing uncertainty.

For the golf swing under SDN, the stochastic extension of the attainable set framework tells us that not all nominally reachable configurations are equally \emph{robustly} reachable. A configuration in the interior of the deterministic attainable set may lie on the boundary of the stochastic attainable set if noise-driven perturbations commonly push trajectories away from it. By contrast, configurations that are "stable" with respect to noise (in the direction of the control Lie algebra) remain deeply inside the stochastic attainable set. The funnel computed in Chapter 7 with stochastic cost weighting thus encodes not just reachability, but \textbf{robust attainment}---guarantee of success under unavoidable noise.

% ══════════════════════════════════════════════
\section{Summary}
% ══════════════════════════════════════════════

\begin{center}
\renewcommand{\arraystretch}{1.4}
\begin{tabular}{@{} l p{7.5cm} @{}}
  \toprule
  \textbf{Concept} & \textbf{Implication for Control} \\
  \midrule
  Signal-Dependent Noise (SDN) & Muscle noise variance proportional to command: $\text{Var}(\tilde{u}) = c u^2$. Requires bounded control effort for repeatability. \\
  Minimum Variance Theory & Explains smooth, controlled velocity profiles as solutions to endpoint error variance minimization under SDN. \\
  Stochastic OCP & Opt